\documentclass{article}

\usepackage{graphicx}
\usepackage{amsmath}
\usepackage{amssymb}
\usepackage{amsthm}
\usepackage{mathtools}
\usepackage[most]{tcolorbox}
\usepackage{bm}
\usepackage{hyperref}

\newcommand{\loss}{\mathcal{L}}

\title{Off-Diagonal Inverse}
\date{\today}

\begin{document}

\maketitle

\section{Problem Statement}
Given a model with $N$ parameters $\theta(t)=\{\theta^i(t)\}_{i=1}^N$ and a loss function 
$\loss:\mathbb{R}^N\to\mathbb{R}$, 
our goal is to find the most optimal parameter update $\delta \theta$ at each time step $t$.
That is, in gradient descent we have
$$\frac{d \theta^i}{dt} = - g^{ij} \frac{\partial \loss}{\partial \theta^j} 
\implies 
\delta \theta^i = - \eta g^{ij} \frac{\partial \loss}{\partial \theta^j}$$
where $\eta\in\mathbb{R}^{>0}$ is the learning rate and $g^{ij}$ is the inverse of the induced metric on the parameter space.

\section{Review of Diagonal Metric Update}

The original ambient metric in the paper is 
$$h = \begin{pmatrix} \gamma & \vec{0} \\ \vec{0}^\top & 1 \end{pmatrix}\in\mathbb{R}^{(N+1)\times(N+1)}$$ 
where $\gamma\in\mathbb{R}^{N\times N}$.
What we really care about is the induced metric on the parameter space, which is given by the pullback of $h$ under the embedding 
$\phi:\mathbb{R}^N\to\mathbb{R}^{N+1}$ defined by 
$\phi(\theta^1,\ldots,\theta^N) = (\theta^1,\ldots,\theta^N,f(\theta^1,\ldots,\theta^N))$ 
where $f:\mathbb{R}^N\to\mathbb{R}$ is some scalar function which we take to be the loss function $\mathcal{L}$.
The induced metric $g$ on the parameter space is given by
$$g_{ij} 
= h_{\alpha\beta} \frac{\partial \phi^\alpha}{\partial \theta^i} \frac{\partial \phi^\beta}{\partial \theta^j} 
= \gamma_{ij} + \frac{\partial \mathcal{L}}{\partial \theta^i} \frac{\partial \mathcal{L}}{\partial \theta^j}.$$
Then, to calculate the parameter updates $\delta \theta^i$ we need $g^{ij}$, the inverse of $g_{ij}$.
We can compute this inverse using the Sherman-Morrison-Woodbury formula:
\begin{tcolorbox}[colback=gray!8, colframe=gray!22, breakable, title=\textbf{Sherman-Morrison-Woodbury Formula}, coltitle=black]
Let $A\in\mathbb{R}^{n\times n}$ be an invertible matrix, and let $U,V\in\mathbb{R}^{n\times k}$ be matrices. 
Then, if $C = A + U V^\top$ is invertible, we have
$$C^{-1} = A^{-1} - A^{-1} U (I_k + V^\top A^{-1} U)^{-1} V^\top A^{-1}.$$
\end{tcolorbox}
\noindent This gives us
$$g^{ij} 
= \gamma^{ij} - \frac{\gamma^{ik}\gamma^{jl} \frac{\partial \mathcal{L}}{\partial \theta^k} \frac{\partial \mathcal{L}}{\partial \theta^l}}{1 + \gamma^{mn} \frac{\partial \mathcal{L}}{\partial \theta^m} \frac{\partial \mathcal{L}}{\partial \theta^n}}
= \gamma^{ij} - \frac{l^i l^j}{1+l_k l^k}$$
where
$$l_i \coloneqq \frac{\partial\mathcal{L}}{\partial \theta^i} \implies l^i = \gamma^{ij} l_j.$$
Thus, the parameter updates are given by
\begin{align*}
    \delta \theta^i =& -\eta \left( \gamma^{ij} - \frac{l^i l^j}{1+l_k l^k} \right) l_j \\
    =& -\eta l^i \left( 1 - \frac{l_j l^j}{1+l_k l^k} \right) \\
    =& -\eta l^i \left( \frac{1}{1+l_k l^k} \right).
\end{align*}

So to summarize, (1) we start with an ambient metric $h$ which induces a metric $g$ on the parameter space, 
(2) we compute the inverse $g^{ij}$, and 
(3) we use $g^{ij}$ to compute the parameter updates $\delta \theta^i$.

\section{Off-Diagonal Metric Update}

Now, we consider a more general ambient metric of the form
$$h = \begin{pmatrix} \gamma & \vec{a} \\ \vec{b}^\top & 1 \end{pmatrix}\in\mathbb{R}^{(N+1)\times(N+1)}$$
where $\vec{a},\vec{b}\in\mathbb{R}^N$.
Intuitively, the vectors $\vec{a}$ and $\vec{b}$ introduce off-diagonal couplings between the parameter space and the loss function dimension.
The induced metric on the parameter space is then given by
\begin{align*}
    g_{ij} 
    &= h_{\alpha\beta} \frac{\partial \phi^\alpha}{\partial \theta^i} \frac{\partial \phi^\beta}{\partial \theta^j} \\
    &= \gamma_{ij} + a_i \frac{\partial \loss}{\partial \theta^j} + b_j \frac{\partial \mathcal{L}}{\partial \theta^i} + \frac{\partial \mathcal{L}}{\partial \theta^i} \frac{\partial \mathcal{L}}{\partial \theta^j} \\
    &= \gamma_{ij} + a_i l_j + b_j l_i + l_i l_j.
\end{align*}
To compute the inverse $g^{ij}$, we notice that $g$ can be expressed as a rank-2 update to $\gamma$, 
that is for $U=[a, l], V=[l, b + l]\in\mathbb{R}^{N\times 2}$ we have
\begin{align*}
\gamma_{ij} + U_{ik} V_{jk} 
&= \gamma_{ij} + a_i l_j + (b_j + l_j) l_i \\
&= \gamma_{ij} + a_i l_j + b_j l_i + l_i l_j \\
&= g_{ij}
\end{align*} and so we can apply the Sherman-Morrison-Woodbury formula again to get
\begin{align*}
g^{ij} 
&= \gamma^{ij} - \gamma^{ik} U_{k\alpha} (\delta_{\alpha\beta} + V_{\beta l} \gamma^{lm} U_{m\alpha})^{-1} V_{\beta n} \gamma^{nj} \\
&= \gamma^{ij} - \gamma^{ik} (a_k, l_k) \begin{pmatrix}
1 + (b_m + l_m) \gamma^{mn} a_n & (b_m + l_m) \gamma^{mn} l_n \\
l_m \gamma^{mn} a_n & 1 + l_m \gamma^{mn} l_n
\end{pmatrix}^{-1} \begin{pmatrix}
b_n + l_n \\ l_n
\end{pmatrix} \gamma^{nj} \\
&= \gamma^{ij} - \gamma^{ik} (a_k, l_k) 
\begin{pmatrix}
1 + b_n a^n + l_n a^n & b_n l^n + l_n l^n \\
l_n a^n & 1 + l_n l^n
\end{pmatrix}^{-1} \begin{pmatrix}
l_n \\ b_n + l_n
\end{pmatrix} \gamma^{nj} \\
&= \gamma^{ij} - \gamma^{ik} \frac{1}{D} (a_k, l_k) 
\begin{pmatrix}
1 + l_n l^n & - (b_n l^n + l_n l^n) \\
- l_n a^n & 1 + b_n a^n + l_n a^n
\end{pmatrix} \begin{pmatrix}
l_n \\ b_n + l_n
\end{pmatrix} \gamma^{nj}
\end{align*} where $D = (1 + b_n a^n + l_n a^n)(1 + l_n l^n) - (b_n l^n + l_n l^n)(l_n a^n)$.
\begin{tcolorbox}[colback=gray!8, colframe=gray!22, breakable, title=\textbf{Quick Sanity Check}, coltitle=black]
If we set $\vec{a}=\vec{0}$ and $\vec{b}=\vec{0}$, we recover the diagonal metric case:
\begin{align*}
g^{ij} 
&= \gamma^{ij} - \gamma^{ik}  
\frac{1}{1 + l_n l^n} (0, l_k) \begin{pmatrix}
1 + l_n l^n & - l_n l^n \\
0 & 1
\end{pmatrix} \begin{pmatrix}
l_n \\ l_n
\end{pmatrix} \gamma^{nj} \\
&= \gamma^{ij} - \gamma^{ik}  
\frac{1}{1 + l_n l^n} (0, l_k) \begin{pmatrix}
(1 + l_n l^n) l_n - (l_n l^n) l_n \\
l_n
\end{pmatrix} \gamma^{nj} \\
&= \gamma^{ij} - \gamma^{ik} 
\frac{1}{1 + l_n l^n} (0, l_k)\begin{pmatrix}
l_n \\ l_n
\end{pmatrix} \gamma^{nj} \\
&= \gamma^{ij} - \frac{\gamma^{ik} l_k l_n \gamma^{nj}}{1 + l_m l^m} \\
&= \gamma^{ij} - \frac{l^i l^j}{1 + l_k l^k}
\end{align*} which matches our previous result.
\end{tcolorbox}
\noindent Thus, the parameter updates are given by
\begin{align*}
    \delta \theta^i 
    =& -\eta g^{ij} l_j \\
    =& -\eta \left[ \gamma^{ij} - \gamma^{ik} \frac{1}{D} (a_k, l_k) 
    \begin{pmatrix}
    1 + l_n l^n & - (b_n l^n + l_n l^n) \\
    - l_n a^n & 1 + b_n a^n + l_n a^n
    \end{pmatrix} \begin{pmatrix}
    l_n \\ b_n + l_n
    \end{pmatrix} \right] l_j \\
    =& -\eta \left[ l^i - \gamma^{ik} \frac{1}{D} (a_k, l_k) 
    \begin{pmatrix}
    1 + l_n l^n & - (b_n l^n + l_n l^n) \\
    - l_n a^n & 1 + b_n a^n + l_n a^n
    \end{pmatrix} \begin{pmatrix}
    l_n \\ b_n + l_n
    \end{pmatrix} \right] \\
    =& -\eta \left[ l^i - \gamma^{ik} \frac{1}{D} \left( a_k \left[ (1 + l_n l^n) l_n - (b_n l^n + l_n l^n)(b_n + l_n) \right] + l_k \left[ - l_n a^n l_n + (1 + b_n a^n + l_n a^n)(b_n + l_n) \right] \right) \right] \\
    =& -\eta \left[ l^i - \frac{1}{D} \left( a^i \left[ (1 + l_n l^n) l_n - (b_n l^n + l_n l^n)(b_n + l_n) \right] + l^i \left[ - l_n a^n l_n + (1 + b_n a^n + l_n a^n)(b_n + l_n) \right] \right) \right]
\end{align*}

\section{Results}


\end{document}
